\section{Fundamentação Teórica}
\label{sec:fundamentacao}

Esta seção apresenta os conceitos de \textit{DevOps}, Infraestrutura como Código, Inteligência Artificial generativa e pipelines de \textit{CI/CD} que sustentam a geração, a validação e a comparação dos artefatos avaliados neste trabalho.

\subsection{\textit{DevOps} e Infraestrutura como Código}

\textit{DevOps} integra desenvolvimento e operações para aumentar a velocidade e a confiabilidade das entregas. Para Kim et al. \cite{kim2016devops}, essa cultura se sustenta nos Três Caminhos (\textit{The Three Ways}): fluxo de trabalho, ciclos de \textit{feedback} e aprendizado contínuo. Em ambientes \textit{cloud}, a automação reduz atividades manuais e permite incorporar alterações de forma mais frequente e controlada.

A Infraestrutura como Código (\textit{IaC}) leva esse princípio à infraestrutura: recursos computacionais passam a ser definidos em arquivos versionáveis, processados por ferramentas como Terraform, AWS CloudFormation e OpenTofu, e submetidos a práticas de engenharia de software como versionamento, revisão de código e automação de mudanças \cite{morris2016infrastructure}. Como qualquer artefato de software, porém, um \textit{template} de \textit{IaC} pode conter erros, permissões excessivas ou divergências entre o estado esperado e o real da infraestrutura, de modo que sua validade sintática não garante, por si só, que o recurso provisionado esteja correto ou seguro. Por isso, revisão e validação automatizada são mecanismos importantes antes da aplicação em ambientes produtivos.

\subsection{Inteligência Artificial Generativa e Modelos de Linguagem}

\textit{Large Language Models} (\textit{LLMs}) são modelos baseados na arquitetura \textit{Transformer} \cite{vaswani2017attention}, capazes de transformar descrições em linguagem natural em código, ainda que sem garantia de precisão \cite{chen2021codex}. Neste trabalho, essa capacidade é usada para gerar arquivos Terraform a partir de requisitos funcionais, e a saída depende do \textit{prompt} e das informações fornecidas, como padrões de nomenclatura, requisitos de segurança e restrições de implementação, que variam entre organizações.

Dois riscos motivam a validação adotada neste experimento. Primeiro, os LLMs podem alucinar, ou seja, produzir informações incorretas de forma aparentemente plausível \cite{ji2023hallucination}. Na geração de código, isso pode resultar em propriedades inexistentes ou atendimento apenas parcial aos requisitos, e diferentes execuções da mesma solicitação podem produzir respostas distintas. Segundo, uma resposta sintaticamente correta não é necessariamente segura: Pearce et al. \cite{pearce2022copilot} identificaram vulnerabilidades em código produzido com apoio de modelos generativos, e em IaC isso pode significar permissões excessivas ou exposição pública sem que o \textit{template} deixe de ser válido. Por isso, neste trabalho, os artefatos gerados sem e com contexto organizacional são submetidos às mesmas verificações aplicadas aos \textit{templates} manuais.

\subsection{Pipelines de Integração e Entrega Contínua}

Pipelines de \textit{CI/CD} automatizam validação, testes e disponibilização de software e infraestrutura, aumentando a repetibilidade das operações e o \textit{feedback} sobre as alterações realizadas \cite{humble2010continuous}. Em \textit{IaC}, isso permite executar de forma padronizada verificações de sintaxe, conformidade e segurança, reduzindo a dependência de checagens manuais. É esse mecanismo que, neste trabalho, permite submeter \textit{templates} manuais e gerados por IA aos mesmos controles, produzindo métricas comparáveis entre os três cenários experimentais.
